\documentclass[11pt,a4paper]{article}
\usepackage[utf8]{inputenc}\usepackage[T1]{fontenc}\usepackage[italian,english]{babel}
\usepackage{amsmath,amssymb,amsthm}\usepackage[margin=2.2cm]{geometry}\usepackage{longtable,booktabs,array,calc}
\usepackage{listings,xcolor}\usepackage{hyperref}\usepackage{url}
\providecommand{\tightlist}{\setlength{\itemsep}{0pt}\setlength{\parskip}{0pt}}
\lstset{basicstyle=\ttfamily\scriptsize,breaklines=true,frame=single,captionpos=b,columns=fullflexible,keepspaces=true,inputencoding=utf8,extendedchars=true}
\title{Proof Pursuit --- Problem 4, part 5\\ \large prove, decisioni degli agenti, codice verificato, letteratura}
\author{Team Proof Pursuit (Researcher: T. Tumini; Referee: gabundos; Creative: M. Renzi; gio3r3gio)}
\date{\today}
\begin{document}\maketitle\tableofcontents
\section*{Come leggere questo rapporto}
Per ogni cella: la richiesta ufficiale, lo stato dichiarato dal team (mai ``risolto'' senza revisione umana), la consegna
con la prova, poi la traccia delle decisioni degli agenti (Researcher: famiglia e motivo; Referee: verdetto ed errore
fatale; Creative: quando e perché è intervenuto) e il codice su cui il tentativo poggia con l'esito della riesecuzione
fidata. DIMOSTRATO, CITATO, VERIFICATO SU INTERVALLO FINITO e CONGETTURA restano distinti. L'architettura del sistema
è descritta nel README del repository.
\section{Problem 4 — Disjoint congruence classes and large gcd}
\subsection{Cella 5 (8 punti) — stato: parziale}
\subparagraph{Parte 5 (C5) --- The certified
boundary}\label{parte-5-c5-the-certified-boundary}

\textbf{Punteggio:} 8 points · \textbf{Valutazione:} Judged

The main certification cell: an initial range of sizes as long as you
can make it, plus two isolated sizes further out. Determine the largest
\(k\) for which you can certify the statement for all sizes up to and
including \(k\), and certify it. The range you claim must be contiguous,
and if your argument at a given size assumes that all smaller sizes have
already been settled you must say so. Then decide the two isolated sizes
\(k = 24\) and \(k = 30\): show that neither is the least size at which
the statement can fail.

For this cell the exhaustiveness certificate of the hand-in rules is not
enough on its own. Add:

\begin{enumerate}
\def\labelenumi{(\roman{enumi})}
\item
  a demonstration that your method is not vacuous. Construct a case in
  which an admissible configuration genuinely exists, run your machinery
  on it unmodified, and show it returns that configuration. A method
  that reports ``nothing survives'' at every size it is pointed at is
  indistinguishable from a method with a bug, and will be graded as one;
\item
  for every object your pruning leaves undecided, at every \(k\) you
  claim --- not a sample --- an individual decision, plus the smallest
  part of that object which already forces the decision, plus a proof
  that no smaller part does;
\item
  the exact list, not merely the count, of what survives at each \(k\);
\item
  every pruning rule you use beyond those you have proved, proved. If
  your search needs a rule you invented to finish a size, that rule is
  part of your claim: state it, prove it, and show the survivor list is
  unchanged when you switch it off. A size that only closes with an
  unproved rule is not certified;
\item
  a second implementation, written independently of your first, that
  differs in method --- not the same algorithm twice --- and the two
  survivor lists compared elementwise at every \(k\) you claim. Report
  any difference rather than reconciling it silently: a disagreement
  means one of them is wrong, and finding which is part of the cell.
\end{enumerate}

\textbf{Enunciato protetto dato agli agenti (state.json).} \texttt{Determine the largest k for which the statement is certified for all sizes up to k (contiguous range), with the strengthened certificate (i)-(v): non-vacuity demonstration, individual decisions for survivors with minimal forcing part, exact survivor lists, all pruning rules proved and tested off, a second independent implementation compared elementwise. Also decide k = 24 and k = 30: show neither is the least size at which the statement can fail.}

\textbf{Evidenza registrata in STATUS.md.} richiede 2 implementazioni indipendenti; bozza parziale in inglese in submission/ (parziali.py)

\subsubsection{Consegna (prova)}
\paragraph{Problema 4 --- Parte 5 --- bozza di consegna
PARZIALE}\label{problema-4-parte-5-bozza-di-consegna-parziale}

\textbf{Stato dichiarato: PARZIALE.} Nessuna soluzione completa: qui
sotto ciò che è stabilito, la formalizzazione, la posizione rispetto
alla letteratura e ciò che resta aperto. Nulla è dichiarato dimostrato
oltre quanto scritto.

\subparagraph{1. Risultato e ambito}\label{risultato-e-ambito}

\textbf{Richiesta ufficiale.} \#\# Parte 5 (C5) --- The certified
boundary

\textbf{Punteggio:} 8 points · \textbf{Valutazione:} Judged

The main certification cell: an initial range of sizes as long as you
can make it, plus two isolated sizes further out. Determine the largest
\(k\) for which you can certify the statement for all sizes up to and
including \(k\), and certify it. The range you claim must be contiguous,
and if your argument at a given size assumes that all smaller sizes have
already been settled you must say so. Then decide the two isolated sizes
\(k = 24\) and \(k = 30\): show that neither is the least size at which
the statement can fail.

For this cell the exhaustiveness certificate of the hand-in rules is not
enough on its own. Add:

\begin{enumerate}
\def\labelenumi{(\roman{enumi})}
\item
  a demonstration that your method is not vacuous. Construct a case in
  which an admissible configuration genuinely exists, run your machinery
  on it unmodified, and show it returns that configuration. A method
  that reports ``nothing survives'' at every size it is pointed at is
  indistinguishable from a method with a bug, and will be graded as one;
\item
  for every object your pruning leaves undecided, at every \(k\) you
  claim --- not a sample --- an individual decision, plus the smallest
  part of that object which already forces the decision, plus a proof
  that no smaller part does;
\item
  the exact list, not merely the count, of what survives at each \(k\);
\item
  every pruning rule you use beyond those you have proved, proved. If
  your search needs a rule you invented to finish a size, that rule is
  part of your claim: state it, prove it, and show the survivor list is
  unchanged when you switch it off. A size that only closes with an
  unproved rule is not certified;
\item
  a second implementation, written independently of your first, that
  differs in method --- not the same algorithm twice --- and the two
  survivor lists compared elementwise at every \(k\) you claim. Report
  any difference rather than reconciling it silently: a disagreement
  means one of them is wrong, and finding which is part of the cell.
\end{enumerate}

\textbf{Cosa consegniamo.} Nessun certificato. Osservazione: per
\(k=24,30\) la forza bruta su \(L_k\) è impraticabile; serve un
argomento di riduzione di taglia (se un controesempio esiste a \(k\), ne
esiste uno a un \(k'<k\)).

\subparagraph{2. Dimostrazione}\label{dimostrazione}

Non sviluppato.

\subparagraph{3. Verifica: istruzioni, dipendenze,
tempi}\label{verifica-istruzioni-dipendenze-tempi}

Codice disponibile (Python 3, libreria standard; ogni script gira in
meno di un minuto): - Nessun codice ancora.

\subparagraph{4. Fonti e contributo}\label{fonti-e-contributo}

Letteratura arXiv (ricerca deterministica
\texttt{tools/cerca\_letteratura.sh}, abstract letti, non usata come
prova): - arXiv:2603.26043v1 --- Finiteness of Disjoint Covering Systems
with Precisely One Repeated Modulus (Yu Hashimoto, 2026); solo abstract
letto. - arXiv:1511.04293v1 --- Searching for Disjoint Covering Systems
with Precisely One Repeated Modulus (Shalosh B. Ekhad, Aviezri S.
Fraenkel, Doron Zeilberger, 2015); solo abstract letto. -
arXiv:2607.24655v1 --- On the problem of large gcd for disjoint residue
classes (Jan Fornal, Yu-Chen Sun, 2026); solo abstract letto. -
arXiv:2608.15873v1 --- Two Questions on \(G\)-harmonic Tuples (Murali
Menon, 2026); solo abstract letto. Come parte 3.

\subparagraph{5. Limiti e parti
irrisolte}\label{limiti-e-parti-irrisolte}

Tutto.

\subsubsection{Decisioni degli agenti}
\paragraph{Run \texttt{p4\_c5}}

\subsubsection{Codice su cui poggia il tentativo}
Nessun codice nei tentativi.

\section{arXiv literature consulted}
\begin{thebibliography}{99}
\bibitem{2603.26043v1} Yu Hashimoto. \emph{Finiteness of Disjoint Covering Systems with Precisely One Repeated Modulus}. arXiv:2603.26043v1 (2026). \url{http://arxiv.org/abs/2603.26043v1}
\bibitem{1511.04293v1} Shalosh B. Ekhad, Aviezri S. Fraenkel, Doron Zeilberger. \emph{Searching for Disjoint Covering Systems with Precisely One Repeated Modulus}. arXiv:1511.04293v1 (2015). \url{http://arxiv.org/abs/1511.04293v1}
\bibitem{2607.24655v1} Jan Fornal, Yu-Chen Sun. \emph{On the problem of large gcd for disjoint residue classes}. arXiv:2607.24655v1 (2026). \url{http://arxiv.org/abs/2607.24655v1}
\bibitem{2608.15873v1} Murali Menon. \emph{Two Questions on \$G\$-harmonic Tuples}. arXiv:2608.15873v1 (2026). \url{http://arxiv.org/abs/2608.15873v1}
\end{thebibliography}
\end{document}
